\newpage
\section{Triển khai chi tiết}

\subsection{Source Code Repository}
 
Toàn bộ mã nguồn được quản lý trên GitHub:
 
\begin{center}
\url{https://github.com/TranQuocViet26701/lsh-book-recommendation}
\end{center}

\subsection{Bộ dữ liệu --- Project Gutenberg}

\textbf{Nguồn:} Project Gutenberg (\url{https://www.gutenberg.org/}) --- thư viện sách miễn phí với hơn 70,000 đầu sách.

\begin{table}[h!]
\centering
\caption{Thông tin bộ dữ liệu Project Gutenberg}
\label{tab:dataset}
\begin{tabularx}{\textwidth}{|l|X|}
\hline
\textbf{Thuộc tính} & \textbf{Chi tiết} \\
\hline
Quy mô dự kiến & 3,000 -- 10,000 đầu sách (subset có thể điều chỉnh để đánh giá scalability) \\
\hline
Định dạng & Plain text (\texttt{.txt}), mã hóa UTF-8 \\
\hline
Ngôn ngữ & Tiếng Anh (có thể mở rộng đa ngữ) \\
\hline
Kích thước trung bình & $\sim$300 KB -- 1 MB / sách \\
\hline
Tổng dung lượng ước tính & $\sim$5 -- 10 GB (cho 10,000 sách) \\
\hline
Metadata đi kèm & Tên sách, tác giả, thể loại, năm xuất bản (từ Gutenberg catalog) \\
\hline
\end{tabularx}
\end{table}

\textbf{Phương pháp thu thập:} Sử dụng Gutenberg mirror hoặc thư viện Python \texttt{gutenberg} / \texttt{gutenbergpy} để tải batch dữ liệu, sau đó upload lên HDFS.

\subsection{Tiền xử lý dữ liệu (Preprocessing Pipeline)}

Toàn bộ quá trình tiền xử lý được thực hiện song song trên Spark:

\begin{table}[h!]
\centering
\caption{Các bước tiền xử lý dữ liệu}
\label{tab:preprocessing}
\begin{tabularx}{\textwidth}{|c|l|X|}
\hline
\textbf{Bước} & \textbf{Phương pháp} & \textbf{Mục đích} \\
\hline
1 & Loại bỏ header/footer (Regex) & Xóa phần ``Project Gutenberg'' boilerplate ở đầu và cuối sách, giữ lại nội dung thuần túy \\
\hline
2 & Chuyển chữ thường (\texttt{lower()}) & Chuẩn hóa case \\
\hline
3 & Loại bỏ ký tự đặc biệt (Regex) & Giữ lại chữ cái, số, khoảng trắng; giảm nhiễu \\
\hline
4 & Tokenization & Tách từ (word-level hoặc character-level tùy chiến lược shingling) \\
\hline
5 & Loại bỏ stopwords & Sử dụng danh sách stopwords tiếng Anh (NLTK hoặc custom); giảm chiều, tăng chất lượng so sánh \\
\hline
6 & Stemming / Lemmatization \textit{(tùy chọn)} & Porter Stemmer hoặc WordNet Lemmatizer; chuẩn hóa biến thể từ \\
\hline
\end{tabularx}
\end{table}

\subsection{Cấu trúc lưu trữ trên HDFS}

Dữ liệu được tổ chức theo cấu trúc thư mục rõ ràng trên HDFS, tách biệt giữa dữ liệu thô, dữ liệu đã xử lý, và các sản phẩm trung gian:

\begin{lstlisting}[style=CStyle, language=bash, caption={Cấu trúc thư mục trên HDFS}, label={lst:hdfs-structure}]
/project-lsh/
|-- raw/                        # Du lieu sach goc
|   |-- book_00001.txt
|   |-- book_00002.txt
|   +-- ...
|-- cleaned/                    # Du lieu sau tien xu ly
|   |-- book_00001_cleaned.txt
|   +-- ...
|-- signatures/                 # MinHash signatures (Parquet)
|   +-- minhash_signatures.parquet
|-- lsh_index/                  # LSH bucket assignments
|   +-- lsh_buckets.parquet
|-- metadata/                   # Metadata sach (Parquet)
|   +-- books_metadata.parquet
+-- results/                    # Ket qua query
    +-- similarity_results.parquet
\end{lstlisting}

\subsection{Thuật toán LSH --- Chi tiết kỹ thuật}

\subsubsection{Shingling (k-gram)}

Mỗi tài liệu được biểu diễn thành một \textbf{tập hợp các k-shingle} (chuỗi con liên tiếp gồm $k$ ký tự hoặc $k$ từ):

\begin{itemize}
    \item \textbf{Character-level shingle} ($k = 5$--$10$): Phù hợp cho phát hiện sao chép, ít nhạy cảm với lỗi chính tả.
    \item \textbf{Word-level shingle} ($k = 2$--$5$): Phù hợp cho so sánh ngữ nghĩa ở mức cụm từ.
    \item \textbf{Lựa chọn mặc định:} Word-level shingle với $k = 3$ (trigram) --- cân bằng giữa granularity và hiệu quả.
\end{itemize}

\subsubsection{MinHash Signatures}

MinHash là kỹ thuật ước lượng Jaccard Similarity giữa hai tập hợp mà không cần so sánh trực tiếp:

\begin{itemize}
    \item Sử dụng $n$ hàm hash ($n = 100$--$200$) để tạo MinHash signature cho mỗi tài liệu.
    \item Mỗi signature là một vector gồm $n$ giá trị nguyên.
    \item \textbf{Tính chất:} Xác suất hai signature trùng tại một vị trí bằng đúng Jaccard Similarity của hai tập shingle gốc:
\end{itemize}

\begin{equation}
    P(h(A) = h(B)) = J(A, B) = \frac{|A \cap B|}{|A \cup B|}
\end{equation}

\subsubsection{LSH Banding Technique}

Chia mỗi MinHash signature (gồm $n$ giá trị) thành $b$ band, mỗi band gồm $r$ row ($n = b \times r$):

\begin{itemize}
    \item Hai tài liệu trở thành \textbf{candidate pair} nếu chúng hash vào cùng bucket ở \textbf{ít nhất 1 band}.
    \item \textbf{S-curve threshold:} Ngưỡng tương đồng xấp xỉ:
\end{itemize}

\begin{equation}
    t \approx \left(\frac{1}{b}\right)^{1/r}
\end{equation}

\textbf{Ví dụ:} Với $n = 100$, $b = 20$, $r = 5$ $\rightarrow$ threshold $\approx 0.55$ (các cặp có Jaccard $\geq 0.55$ sẽ có xác suất cao trở thành candidate).

\begin{figure}[h!]
\centering
\begin{tikzpicture}[scale=0.85]
    \draw[->] (0,0) -- (8,0) node[right] {\small Jaccard Similarity $s$};
    \draw[->] (0,0) -- (0,5.5) node[above] {\small $P(\text{candidate})$};
    
    % S-curve
    \draw[blue, thick, domain=0:7.5, samples=100] plot (\x, {5*(1-(1-(\x/7.5)^5)^20)});
    
    % Threshold line
    \draw[red, dashed] (4.1,0) -- (4.1,5) node[above, font=\small] {$t \approx 0.55$};
    
    % Step function (ideal)
    \draw[gray, dashed, thick] (0,0) -- (4.1,0) -- (4.1,5) -- (7.5,5);
    
    % Labels
    \node[font=\scriptsize] at (2, 4.5) {S-curve (LSH)};
    \node[font=\scriptsize, gray] at (6, 1) {Ideal};
    
    % Ticks
    \foreach \x/\label in {0/0, 1.875/0.25, 3.75/0.5, 5.625/0.75, 7.5/1.0}
        \draw (\x, -0.1) -- (\x, 0.1) node[below=2pt, font=\scriptsize] {\label};
    \foreach \y/\label in {0/0, 2.5/0.5, 5/1.0}
        \draw (-0.1, \y) -- (0.1, \y) node[left=2pt, font=\scriptsize] {\label};
\end{tikzpicture}
\caption{S-curve của LSH: xác suất trở thành candidate pair theo Jaccard Similarity ($b=20$, $r=5$)}
\label{fig:scurve}
\end{figure}

\subsubsection{Hàm đo tương đồng}

\begin{table}[h!]
\centering
\caption{So sánh các hàm đo tương đồng}
\label{tab:similarity-measures}
\begin{tabularx}{\textwidth}{|l|l|X|}
\hline
\textbf{Hàm đo} & \textbf{Công thức} & \textbf{Phù hợp cho} \\
\hline
Jaccard Similarity \textit{(chính)} & $J(A,B) = \frac{|A \cap B|}{|A \cup B|}$ & So sánh tập hợp shingle --- mặc định cho LSH MinHash \\
\hline
Cosine Similarity \textit{(tham khảo)} & $\cos(A,B) = \frac{A \cdot B}{\|A\| \cdot \|B\|}$ & So sánh vector TF-IDF (cần LSH variant: SimHash) \\
\hline
Hamming Distance \textit{(tham khảo)} & Số vị trí bit khác nhau & So sánh binary fingerprint \\
\hline
\end{tabularx}
\end{table}

Đề tài sử dụng \textbf{Jaccard Similarity} làm hàm đo chính vì nó tự nhiên phù hợp với biểu diễn shingle-set và được hỗ trợ trực tiếp bởi MinHash.

\subsection{Đánh giá LSH trong bài toán Similarity Join (MRS Join)}

\subsubsection{LSH phù hợp với Similarity Join ở mức độ nào?}

\begin{table}[h!]
\centering
\caption{Đánh giá mức độ phù hợp của LSH cho Similarity Join}
\label{tab:lsh-evaluation}
\begin{tabularx}{\textwidth}{|l|X|}
\hline
\textbf{Tiêu chí} & \textbf{Đánh giá} \\
\hline
Tốc độ & LSH giảm số cặp cần so sánh từ $O(n^2)$ xuống $O(n \cdot k)$ với $k \ll n$, rất hiệu quả cho tập dữ liệu lớn \\
\hline
Chất lượng & LSH là phương pháp xấp xỉ --- có thể bỏ sót một số cặp tương tự (false negatives) hoặc trả về cặp không thực sự tương tự (false positives). Chất lượng phụ thuộc vào việc chọn tham số $b$, $r$ \\
\hline
Khả năng phân tán & Rất phù hợp với MapReduce/Spark: mỗi band có thể xử lý song song, bucketing tương tự reduce-by-key \\
\hline
Trade-off & Tăng $b$ (nhiều band) $\rightarrow$ tăng recall nhưng tăng false positives. Tăng $r$ (nhiều row/band) $\rightarrow$ tăng precision nhưng giảm recall \\
\hline
\end{tabularx}
\end{table}

\subsubsection{Ứng dụng thực tế của LSH Similarity Join}

\begin{itemize}
    \item \textbf{Near-duplicate detection:} Phát hiện bài viết/tài liệu sao chép trong kho dữ liệu lớn (ứng dụng trong anti-plagiarism, web deduplication).
    \item \textbf{Recommendation systems:} Gợi ý sản phẩm/nội dung tương tự dựa trên đặc trưng nội dung (bài toán của đề tài này).
    \item \textbf{Entity resolution:} Ghép nối các bản ghi trùng lặp trong cơ sở dữ liệu lớn (ví dụ: ghép tên khách hàng, địa chỉ).
    \item \textbf{Genome/protein similarity:} So sánh chuỗi DNA/protein trong tin sinh học.
    \item \textbf{Image retrieval:} Tìm ảnh tương tự trong kho ảnh lớn (kết hợp LSH với feature vector từ CNN).
\end{itemize}
